% Source: https://artofproblemsolving.com/community/c3h1907605_how_to_find_the_answer_to_this_sort_of_question
% Stable identity: aops-topic-1907605-q01
\textbf{Remainders modulo $7$.} Since $2^3=8\equiv1\pmod7$ and $2019=3\cdot673$, we get $2^{2019}=(2^3)^{673}\equiv1$. Since $3^3=27\equiv-1\pmod7$, we get $3^{2019}=(3^3)^{673}\equiv(-1)^{673}=-1$, and $3^{2018}=(3^3)^{672}\cdot3^2\equiv9\equiv2$.

(a) $2^{2019}+3^{2018}\equiv1+2=3\pmod7$, so the remainder is $3$ and the quotient by $7$ is not an integer.

(b) $N\equiv1+(-1)=0\pmod7$, so $7\mid N$.

\textbf{Units digit of $N$.} Units digits of powers of $2$ repeat with period $4$ as $2,4,8,6$, and those of powers of $3$ as $3,9,7,1$. Since $2019\equiv3\pmod4$, $2^{2019}$ ends in $8$ and $3^{2019}$ ends in $7$, so $N$ ends in $5$ ($8+7=15$).

\textbf{Units digit of $\dfrac {N}{7}$.} Let $M=\dfrac {N}{7}$, so $7M=N$ ends in $5$. The units digit of $7M$ depends only on the units digit $u$ of $M$: $7u$ must end in $5$. Among $u=0,1,\dots,9$, only $u=5$ works ($7\cdot5=35$). So $\dfrac {N}{7}$ ends in $\boxed{5}$.

(In general, to find the last digit of $\dfrac {N}{7}$ for an integer $N$ divisible by $7$, find the last digit $d$ of $N$ and solve $7u\equiv d\pmod{10}$; since $7\cdot3=21\equiv1$, $u\equiv3d\pmod{10}$.)
